\begin{theorem}[Periodic holomorphic function on the upper half-plane descends to the punctured disk and has a Fourier--Laurent expansion]
\label{theorem:periodic_holomorphic_function_on_upper_half_plane_descends_to_punctured_disk_and_has_fourier_laurent_expansion}
\TODO{AI generated, verify}
Let $h \in \mathbb{R}_{>0}$, let $\mathbb{H}$ denote the \CrefAndHyperrefIfExist{definition:upper_half_plane_in_the_complex_plane}{upper half-plane}, and let $f : \mathbb{H} \to \mathbb{C}$ be a \CrefAndHyperrefIfExist{definition:holomorphic_function_from_a_subset_of_Cn_to_Cm_and_derivative_of_a_holomorphic_function_at_a_point}{holomorphic function} satisfying
$$f(z + h) = f(z) \quad \text{for all } z \in \mathbb{H}.$$

Define the map $\phi : \mathbb{H} \to \mathbb{C}$ by
$$\phi(z) = e^{2\pi i z / h},$$
where $e^w$ denotes the \CrefAndHyperrefIfExist{definition:complex_exponential_function}{complex exponential function}. Then:

\begin{enumerate}
  \item The image of $\phi$ is the \CrefAndHyperrefIfExist{definition:punctured_disk_in_the_complex_plane}{punctured unit disk}
  $$D^* = \{q \in \mathbb{C} : 0 < |q| < 1\},$$
  and $\phi : \mathbb{H} \to D^*$ is a local \CrefAndHyperrefIfExist{definition:biholomorphic_map_of_open_subsets_of_general_dimensional_complex_space}{biholomorphism} whose fibres are precisely the orbits of the translation $z \mapsto z + h$\CrefIfExists{definition:holomorphic_function_from_a_subset_of_Cn_to_Cm_and_derivative_of_a_holomorphic_function_at_a_point}.

  \item There exists a unique holomorphic function $g : D^* \to \mathbb{C}$ such that
  $$f(z) = g(e^{2\pi i z / h}) \quad \text{for all } z \in \mathbb{H}.$$

  \item The function $g$ admits a \CrefAndHyperrefIfExist{theorem:laurent_series_expansion_for_holomorphic_functions_on_annuli_in_the_complex_plane}{Laurent expansion}
  $$g(q) = \sum_{n=-\infty}^{\infty} a_n q^n$$
  valid for all $q \in D^*$, which converges absolutely and uniformly on every \CrefAndHyperrefIfExist{definition:annulus_in_the_complex_plane}{annulus} $0 < r_1 \leq |q| \leq r_2 < 1$. The coefficients $a_n$ are given by the contour integral
  $$a_n = \frac{1}{2\pi i} \oint_{|q| = \rho} \frac{g(q)}{q^{n+1}} \, dq$$
  for any $\rho \in (0, 1)$.

  \item Substituting $q = e^{2\pi i z / h}$, the Laurent expansion of $g$ becomes the \CrefAndHyperrefIfExist{definition:fourier_expansion_of_a_periodic_holomorphic_function_on_the_upper_half_plane}{Fourier expansion of $f$}
  $$f(z) = \sum_{n=-\infty}^{\infty} a_n \, e^{2\pi i n z / h},$$
  converging absolutely and uniformly on every compact subset of $\mathbb{H}$. The coefficients $a_n$ coincide with the Fourier coefficients defined by the contour integral formula
  $$a_n = \frac{1}{h} \int_{z_0}^{z_0 + h} f(z) \, e^{-2\pi i n z / h} \, dz.$$
\end{enumerate}
\end{theorem}

\begin{proof}
\TODO{AI generated, verify}
We prove each part in order.

\begin{enumerate}
  \item For $z = x + iy \in \mathbb{H}$ with $y > 0$, we have
  $$|\phi(z)| = |e^{2\pi i (x + iy) / h}| = e^{-2\pi y / h} < 1,$$
  so $\phi(z) \in D^*$. Conversely, for any $q \in D^*$, write $q = re^{i\theta}$ with $0 < r < 1$ and $\theta \in \mathbb{R}$. Then $z = \frac{h}{2\pi}(\theta - i \log r)$ satisfies $\phi(z) = q$, and $\Im(z) = -\frac{h}{2\pi} \log r > 0$, so $z \in \mathbb{H}$. Hence $\phi(\mathbb{H}) = D^*$.

  The derivative $\phi'(z) = \frac{2\pi i}{h} e^{2\pi i z / h}$ never vanishes on $\mathbb{H}$ because the exponential function never vanishes. By \Cref{proposition:holomorphic_functions_on_open_subsets_of_C_are_continuous} and the inverse function theorem, $\phi$ is a local biholomorphism. Moreover, $\phi(z_1) = \phi(z_2)$ if and only if $e^{2\pi i (z_1 - z_2)/h} = 1$, which holds precisely when $z_1 - z_2 \in h\mathbb{Z}$. Thus the fibres of $\phi$ are the orbits of the translation $z \mapsto z + h$.

  \item Since $f(z + h) = f(z)$, the function $f$ is constant on each fibre of $\phi$. Because $\phi$ is a local biholomorphism, the universal property of the quotient by a local biholomorphism implies that there exists a unique function $g : D^* \to \mathbb{C}$ such that $f = g \circ \phi$. To see that $g$ is holomorphic, fix $q_0 \in D^*$ and choose $z_0 \in \phi^{-1}(q_0)$. Let $U$ be a neighbourhood of $z_0$ on which $\phi$ restricts to a biholomorphism $\phi|_U : U \xrightarrow{\sim} V$. Then $g|_V = f \circ (\phi|_U)^{-1}$, which is a composition of holomorphic maps and hence holomorphic on $V$. Thus $g$ is holomorphic on $D^*$.

  \item The function $g$ is holomorphic on the punctured disk $D^* = D(0, 1) \setminus \{0\}$. By \Cref{theorem:laurent_series_expansion_for_holomorphic_functions_on_annuli_in_the_complex_plane}, $g$ admits a Laurent expansion
  $$g(q) = \sum_{n=-\infty}^{\infty} a_n q^n$$
  valid for all $q \in D^*$, converging absolutely and uniformly on every closed annulus $0 < r_1 \leq |q| \leq r_2 < 1$. The coefficients are given by the contour integral formula.

  \item Substituting $q = e^{2\pi i z / h}$ yields the Fourier expansion of $f$. The convergence properties follow from the convergence of the Laurent expansion of $g$: for any compact subset $K \subset \mathbb{H}$, the image $\phi(K)$ is contained in a closed annulus $r_1 \leq |q| \leq r_2$ with $0 < r_1 < r_2 < 1$ (since $\Im(z)$ is bounded away from $0$ and $\infty$ on $K$), and the Laurent series converges uniformly on that annulus.

  It remains to verify the coefficient formula. Applying the Laurent coefficient formula to $g$ and pulling back the contour integral via $\phi$, we have for any $\rho \in (0, 1)$:
  \begin{align*}
    a_n &= \frac{1}{2\pi i} \oint_{|q| = \rho} \frac{g(q)}{q^{n+1}} \, dq \\
        &= \frac{1}{2\pi i} \int_{z_0}^{z_0 + h} \frac{g(e^{2\pi i z / h})}{e^{2\pi i (n+1) z / h}} \cdot \frac{2\pi i}{h} e^{2\pi i z / h} \, dz \\
        &= \frac{1}{h} \int_{z_0}^{z_0 + h} f(z) \, e^{-2\pi i n z / h} \, dz,
  \end{align*}
  where the path is the image under $\phi^{-1}$ of the circle $|q| = \rho$, which is a horizontal segment of length $h$ in $\mathbb{H}$. The independence of the choice of path and of $z_0$ follows from \Cref{theorem:cauchy_integral_theorem_stating_that_the_contour_integral_of_a_holomorphic_function_along_any_closed_piecewise_continuously_differentiable_path_in_a_simply_connected_domain_vanishes}. Thus the coefficients $a_n$ coincide with the Fourier coefficients defined in \Cref{definition:fourier_expansion_of_a_periodic_holomorphic_function_on_the_upper_half_plane}.
\end{enumerate}
\end{proof}